
\documentclass{article}
\usepackage{colortbl}
\usepackage{makecell}
\usepackage{multirow}
\usepackage{supertabular}

\begin{document}

\newcounter{utterance}

\centering \large Interaction Transcript for game `matchit\_ascii', experiment `similar\_grid\_1', episode 3 with Qwen3.6{-}35B{-}A3B{-}FP8{-}without{-}reasoning.
\vspace{24pt}

{ \footnotesize  \setcounter{utterance}{1}
\setlength{\tabcolsep}{0pt}
\begin{supertabular}{c@{$\;$}|p{.15\linewidth}@{}p{.15\linewidth}p{.15\linewidth}p{.15\linewidth}p{.15\linewidth}p{.15\linewidth}}
    \# & $\;$A & \multicolumn{4}{c}{Game Master} & $\;\:$B\\
    \hline

    \theutterance \stepcounter{utterance}  
    & & \multicolumn{4}{p{0.6\linewidth}}{
        \cellcolor[rgb]{0.9,0.9,0.9}{
            \makecell[{{p{\linewidth}}}]{
                \texttt{\tiny{[P1$\langle$GM]}}
                \texttt{Zúčastňujete sa spoločnej hry na hádanie.} \\
\texttt{Cieľom je zistiť, či je táto mriežka a iný obrázok, ktorý vidím len ja, rovnaké.} \\
\texttt{Najskôr opíšte svoju mriežku. Potom poskytnem svoj opis a môžeme si navzájom klásť otázky o mriežkach, aby sme zistili, či sú rovnaké. Teraz začnite svoj krátky opis mriežky s "POPIS:", za ktorým nasleduje opis. Nepridávajte nič iné. Toto je vaša mriežka:} \\
\texttt{X X X X X} \\
\texttt{X ▢ ▢ ▢ ▢} \\
\texttt{X X X ▢ ▢} \\
\texttt{X ▢ ▢ ▢ ▢} \\
\texttt{X ▢ ▢ ▢ ▢} \\
            }
        }
    }
    & \\ \\

    \theutterance \stepcounter{utterance}  
    & \multicolumn{4}{p{0.6\linewidth}}{
        \cellcolor[rgb]{0.9,0.9,0.9}{
            \makecell[{{p{\linewidth}}}]{
                \texttt{\tiny{[P1$\rangle$GM]}}
                \texttt{POPIS: Mriežka 5x5 pozostávajúca zo znakov X a prázdnych štvorcov ▢. Prvý stĺpec je úplne vyplnený znakmi X. Druhý riadok má na začiatku X a nasleduje štyri prázdne štvorce. Tretí riadok má na začiatku tri X a potom dva prázdne štvorce. Štvrtý a piaty riadok majú na začiatku X a nasledujú štyri prázdne štvorce.} \\
            }
        }
    }
    & & \\ \\

    \theutterance \stepcounter{utterance}  
    & & & \multicolumn{2}{p{0.3\linewidth}}{
        \cellcolor[rgb]{0.9,0.9,0.9}{
            \makecell[{{p{\linewidth}}}]{
                \texttt{\tiny{[GM$|$GM]}}
                \texttt{continue} \\
            }
        }
    }
    & & \\ \\

    \theutterance \stepcounter{utterance}  
    & & \multicolumn{4}{p{0.6\linewidth}}{
        \cellcolor[rgb]{0.9,0.9,0.9}{
            \makecell[{{p{\linewidth}}}]{
                \texttt{\tiny{[GM$\rangle$P2]}}
                \texttt{Zúčastňujete sa spoločnej hry na hádanie.} \\
\texttt{Cieľom je zistiť, či je táto mriežka a iný obrázok, ktorý vidím len ja, rovnaké.} \\
\texttt{Najskôr opíšte svoju mriežku. Potom poskytnem svoj opis a môžeme si navzájom klásť otázky o mriežkach, aby sme zistili, či sú rovnaké. Teraz začnite svoj krátky opis mriežky s "POPIS:", za ktorým nasleduje opis. Nepridávajte nič iné. Toto je vaša mriežka:} \\
\texttt{X ▢ ▢ ▢ ▢} \\
\texttt{X ▢ ▢ ▢ ▢} \\
\texttt{X X X ▢ ▢} \\
\texttt{X ▢ ▢ ▢ ▢} \\
\texttt{X X X X X} \\
            }
        }
    }
    & \\ \\

    \theutterance \stepcounter{utterance}  
    & & & \multicolumn{4}{p{0.6\linewidth}}{
        \cellcolor[rgb]{0.9,0.9,0.9}{
            \makecell[{{p{\linewidth}}}]{
                \texttt{\tiny{[GM$\langle$P2]}}
                \texttt{POPIS: Mriežka 5x5. Prvý stĺpec je plne vyplnený znakom X. Tretí riadok má tri susedné X na začiatku (stĺpce 1{-}3). Piaty riadok je plne vyplnený znakom X. Zvyšné políčka sú prázdne (▢).} \\
            }
        }
    }
     \\ \\

    \theutterance \stepcounter{utterance}  
    & & & \multicolumn{2}{p{0.3\linewidth}}{
        \cellcolor[rgb]{0.9,0.9,0.9}{
            \makecell[{{p{\linewidth}}}]{
                \texttt{\tiny{[GM$|$GM]}}
                \texttt{continue} \\
            }
        }
    }
    & & \\ \\

    \theutterance \stepcounter{utterance}  
    & & \multicolumn{4}{p{0.6\linewidth}}{
        \cellcolor[rgb]{0.9,0.9,0.9}{
            \makecell[{{p{\linewidth}}}]{
                \texttt{\tiny{[GM$\rangle$P2]}}
                \texttt{Toto je môj POPISPOPIS: Mriežka 5x5 pozostávajúca zo znakov X a prázdnych štvorcov ▢. Prvý stĺpec je úplne vyplnený znakmi X. Druhý riadok má na začiatku X a nasleduje štyri prázdne štvorce. Tretí riadok má na začiatku tri X a potom dva prázdne štvorce. Štvrtý a piaty riadok majú na začiatku X a nasledujú štyri prázdne štvorce.} \\
\texttt{Teraz položte otázku, aby ste zistili nové aspekty mojej mriežky, ktoré sa môžu líšiť od vašej mriežky. Začnite s "OTÁZKA:" a nepridávajte nič iné.} \\
            }
        }
    }
    & \\ \\

    \theutterance \stepcounter{utterance}  
    & & & \multicolumn{4}{p{0.6\linewidth}}{
        \cellcolor[rgb]{0.9,0.9,0.9}{
            \makecell[{{p{\linewidth}}}]{
                \texttt{\tiny{[GM$\langle$P2]}}
                \texttt{OTÁZKA: V treťom riadku sú na začiatku tri X, ale v štvrtom riadku je podľa vášho popisu len jedno X na začiatku. Potvrdzujete, že v štvrtom riadku sú všetky ostatné políčka (stĺpce 2{-}5) prázdne?} \\
            }
        }
    }
     \\ \\

    \theutterance \stepcounter{utterance}  
    & & & \multicolumn{2}{p{0.3\linewidth}}{
        \cellcolor[rgb]{0.9,0.9,0.9}{
            \makecell[{{p{\linewidth}}}]{
                \texttt{\tiny{[GM$|$GM]}}
                \texttt{continue} \\
            }
        }
    }
    & & \\ \\

    \theutterance \stepcounter{utterance}  
    & & \multicolumn{4}{p{0.6\linewidth}}{
        \cellcolor[rgb]{0.9,0.9,0.9}{
            \makecell[{{p{\linewidth}}}]{
                \texttt{\tiny{[P1$\langle$GM]}}
                \texttt{Toto je môj POPISPOPIS: Mriežka 5x5. Prvý stĺpec je plne vyplnený znakom X. Tretí riadok má tri susedné X na začiatku (stĺpce 1{-}3). Piaty riadok je plne vyplnený znakom X. Zvyšné políčka sú prázdne (▢).} \\
\texttt{OTÁZKA: V treťom riadku sú na začiatku tri X, ale v štvrtom riadku je podľa vášho popisu len jedno X na začiatku. Potvrdzujete, že v štvrtom riadku sú všetky ostatné políčka (stĺpce 2{-}5) prázdne?Svoju odpoveď začnite s "ODPOVEĎ:" a nepridávajte nič iné.} \\
            }
        }
    }
    & \\ \\

    \theutterance \stepcounter{utterance}  
    & \multicolumn{4}{p{0.6\linewidth}}{
        \cellcolor[rgb]{0.9,0.9,0.9}{
            \makecell[{{p{\linewidth}}}]{
                \texttt{\tiny{[P1$\rangle$GM]}}
                \texttt{ODPOVEĎ: Áno, potvrdzujem, že v štvrtom riadku je na začiatku jedno X a zvyšné políčka (stĺpce 2{-}5) sú prázdne.} \\
            }
        }
    }
    & & \\ \\

    \theutterance \stepcounter{utterance}  
    & & & \multicolumn{2}{p{0.3\linewidth}}{
        \cellcolor[rgb]{0.9,0.9,0.9}{
            \makecell[{{p{\linewidth}}}]{
                \texttt{\tiny{[GM$|$GM]}}
                \texttt{continue} \\
            }
        }
    }
    & & \\ \\

    \theutterance \stepcounter{utterance}  
    & & \multicolumn{4}{p{0.6\linewidth}}{
        \cellcolor[rgb]{0.9,0.9,0.9}{
            \makecell[{{p{\linewidth}}}]{
                \texttt{\tiny{[P1$\langle$GM]}}
                \texttt{Teraz položte otázku, aby ste zistili nové aspekty mojej mriežky, ktoré sa môžu líšiť od vašej mriežky. Začnite s "OTÁZKA:" a nepridávajte nič iné.} \\
            }
        }
    }
    & \\ \\

    \theutterance \stepcounter{utterance}  
    & \multicolumn{4}{p{0.6\linewidth}}{
        \cellcolor[rgb]{0.9,0.9,0.9}{
            \makecell[{{p{\linewidth}}}]{
                \texttt{\tiny{[P1$\rangle$GM]}}
                \texttt{OTÁZKA: V vašom popise uvádzate, že piaty riadok je plne vyplnený znakom X. V mojej mriežke je piaty riadok "X ▢ ▢ ▢ ▢" (len prvé políčko je X). Potvrdzujete, že v piatom riadku sú stĺpce 2 až 5 naozaj vyplnené znakom X?} \\
            }
        }
    }
    & & \\ \\

    \theutterance \stepcounter{utterance}  
    & & & \multicolumn{2}{p{0.3\linewidth}}{
        \cellcolor[rgb]{0.9,0.9,0.9}{
            \makecell[{{p{\linewidth}}}]{
                \texttt{\tiny{[GM$|$GM]}}
                \texttt{continue} \\
            }
        }
    }
    & & \\ \\

    \theutterance \stepcounter{utterance}  
    & & \multicolumn{4}{p{0.6\linewidth}}{
        \cellcolor[rgb]{0.9,0.9,0.9}{
            \makecell[{{p{\linewidth}}}]{
                \texttt{\tiny{[GM$\rangle$P2]}}
                \texttt{ODPOVEĎ: Áno, potvrdzujem, že v štvrtom riadku je na začiatku jedno X a zvyšné políčka (stĺpce 2{-}5) sú prázdne.} \\
\texttt{OTÁZKA: V vašom popise uvádzate, že piaty riadok je plne vyplnený znakom X. V mojej mriežke je piaty riadok "X ▢ ▢ ▢ ▢" (len prvé políčko je X). Potvrdzujete, že v piatom riadku sú stĺpce 2 až 5 naozaj vyplnené znakom X?Svoju odpoveď začnite s "ODPOVEĎ:" a nepridávajte nič iné.} \\
            }
        }
    }
    & \\ \\

    \theutterance \stepcounter{utterance}  
    & & & \multicolumn{4}{p{0.6\linewidth}}{
        \cellcolor[rgb]{0.9,0.9,0.9}{
            \makecell[{{p{\linewidth}}}]{
                \texttt{\tiny{[GM$\langle$P2]}}
                \texttt{ODPOVEĎ: Áno, potvrdzujem, že v piatom riadku sú všetky políčka (stĺpce 1 až 5) vyplnené znakom X.} \\
            }
        }
    }
     \\ \\

    \theutterance \stepcounter{utterance}  
    & & & \multicolumn{2}{p{0.3\linewidth}}{
        \cellcolor[rgb]{0.9,0.9,0.9}{
            \makecell[{{p{\linewidth}}}]{
                \texttt{\tiny{[GM$|$GM]}}
                \texttt{continue} \\
            }
        }
    }
    & & \\ \\

    \theutterance \stepcounter{utterance}  
    & & \multicolumn{4}{p{0.6\linewidth}}{
        \cellcolor[rgb]{0.9,0.9,0.9}{
            \makecell[{{p{\linewidth}}}]{
                \texttt{\tiny{[GM$\rangle$P2]}}
                \texttt{Teraz položte otázku, aby ste zistili nové aspekty mojej mriežky, ktoré sa môžu líšiť od vašej mriežky. Začnite s "OTÁZKA:" a nepridávajte nič iné.} \\
            }
        }
    }
    & \\ \\

    \theutterance \stepcounter{utterance}  
    & & & \multicolumn{4}{p{0.6\linewidth}}{
        \cellcolor[rgb]{0.9,0.9,0.9}{
            \makecell[{{p{\linewidth}}}]{
                \texttt{\tiny{[GM$\langle$P2]}}
                \texttt{OTÁZKA: V druhom riadku je podľa vášho počiatočného popisu len jedno X na začiatku. Potvrdzujete, že v druhom riadku sú stĺpce 2 až 5 prázdne?} \\
            }
        }
    }
     \\ \\

    \theutterance \stepcounter{utterance}  
    & & & \multicolumn{2}{p{0.3\linewidth}}{
        \cellcolor[rgb]{0.9,0.9,0.9}{
            \makecell[{{p{\linewidth}}}]{
                \texttt{\tiny{[GM$|$GM]}}
                \texttt{continue} \\
            }
        }
    }
    & & \\ \\

    \theutterance \stepcounter{utterance}  
    & & \multicolumn{4}{p{0.6\linewidth}}{
        \cellcolor[rgb]{0.9,0.9,0.9}{
            \makecell[{{p{\linewidth}}}]{
                \texttt{\tiny{[P1$\langle$GM]}}
                \texttt{ODPOVEĎ: Áno, potvrdzujem, že v piatom riadku sú všetky políčka (stĺpce 1 až 5) vyplnené znakom X.} \\
\texttt{OTÁZKA: V druhom riadku je podľa vášho počiatočného popisu len jedno X na začiatku. Potvrdzujete, že v druhom riadku sú stĺpce 2 až 5 prázdne?Svoju odpoveď začnite s "ODPOVEĎ:" a nepridávajte nič iné.} \\
            }
        }
    }
    & \\ \\

    \theutterance \stepcounter{utterance}  
    & \multicolumn{4}{p{0.6\linewidth}}{
        \cellcolor[rgb]{0.9,0.9,0.9}{
            \makecell[{{p{\linewidth}}}]{
                \texttt{\tiny{[P1$\rangle$GM]}}
                \texttt{ODPOVEĎ: Áno, potvrdzujem, že v druhom riadku je na začiatku jedno X a stĺpce 2 až 5 sú prázdne.} \\
            }
        }
    }
    & & \\ \\

    \theutterance \stepcounter{utterance}  
    & & & \multicolumn{2}{p{0.3\linewidth}}{
        \cellcolor[rgb]{0.9,0.9,0.9}{
            \makecell[{{p{\linewidth}}}]{
                \texttt{\tiny{[GM$|$GM]}}
                \texttt{continue} \\
            }
        }
    }
    & & \\ \\

    \theutterance \stepcounter{utterance}  
    & & \multicolumn{4}{p{0.6\linewidth}}{
        \cellcolor[rgb]{0.9,0.9,0.9}{
            \makecell[{{p{\linewidth}}}]{
                \texttt{\tiny{[P1$\langle$GM]}}
                \texttt{Teraz položte otázku, aby ste zistili nové aspekty mojej mriežky, ktoré sa môžu líšiť od vašej mriežky. Začnite s "OTÁZKA:" a nepridávajte nič iné.} \\
            }
        }
    }
    & \\ \\

    \theutterance \stepcounter{utterance}  
    & \multicolumn{4}{p{0.6\linewidth}}{
        \cellcolor[rgb]{0.9,0.9,0.9}{
            \makecell[{{p{\linewidth}}}]{
                \texttt{\tiny{[P1$\rangle$GM]}}
                \texttt{OTÁZKA: V treťom riadku je podľa vášho popisu tri susedné X na začiatku. Potvrdzujete, že v treťom riadku sú stĺpce 4 a 5 prázdne?} \\
            }
        }
    }
    & & \\ \\

    \theutterance \stepcounter{utterance}  
    & & & \multicolumn{2}{p{0.3\linewidth}}{
        \cellcolor[rgb]{0.9,0.9,0.9}{
            \makecell[{{p{\linewidth}}}]{
                \texttt{\tiny{[GM$|$GM]}}
                \texttt{continue} \\
            }
        }
    }
    & & \\ \\

    \theutterance \stepcounter{utterance}  
    & & \multicolumn{4}{p{0.6\linewidth}}{
        \cellcolor[rgb]{0.9,0.9,0.9}{
            \makecell[{{p{\linewidth}}}]{
                \texttt{\tiny{[GM$\rangle$P2]}}
                \texttt{ODPOVEĎ: Áno, potvrdzujem, že v druhom riadku je na začiatku jedno X a stĺpce 2 až 5 sú prázdne.} \\
\texttt{OTÁZKA: V treťom riadku je podľa vášho popisu tri susedné X na začiatku. Potvrdzujete, že v treťom riadku sú stĺpce 4 a 5 prázdne?Svoju odpoveď začnite s "ODPOVEĎ:" a nepridávajte nič iné.} \\
            }
        }
    }
    & \\ \\

    \theutterance \stepcounter{utterance}  
    & & & \multicolumn{4}{p{0.6\linewidth}}{
        \cellcolor[rgb]{0.9,0.9,0.9}{
            \makecell[{{p{\linewidth}}}]{
                \texttt{\tiny{[GM$\langle$P2]}}
                \texttt{ODPOVEĎ: Áno, potvrdzujem, že v treťom riadku sú stĺpce 4 a 5 prázdne.} \\
            }
        }
    }
     \\ \\

    \theutterance \stepcounter{utterance}  
    & & & \multicolumn{2}{p{0.3\linewidth}}{
        \cellcolor[rgb]{0.9,0.9,0.9}{
            \makecell[{{p{\linewidth}}}]{
                \texttt{\tiny{[GM$|$GM]}}
                \texttt{continue} \\
            }
        }
    }
    & & \\ \\

    \theutterance \stepcounter{utterance}  
    & & \multicolumn{4}{p{0.6\linewidth}}{
        \cellcolor[rgb]{0.9,0.9,0.9}{
            \makecell[{{p{\linewidth}}}]{
                \texttt{\tiny{[GM$\rangle$P2]}}
                \texttt{Teraz položte otázku, aby ste zistili nové aspekty mojej mriežky, ktoré sa môžu líšiť od vašej mriežky. Začnite s "OTÁZKA:" a nepridávajte nič iné.} \\
            }
        }
    }
    & \\ \\

    \theutterance \stepcounter{utterance}  
    & & & \multicolumn{4}{p{0.6\linewidth}}{
        \cellcolor[rgb]{0.9,0.9,0.9}{
            \makecell[{{p{\linewidth}}}]{
                \texttt{\tiny{[GM$\langle$P2]}}
                \texttt{OTÁZKA: V prvom stĺpci sú podľa vášho popisu vo všetkých riadkoch znaky X. Potvrdzujete, že v prvom stĺpci sú naozaj znaky X vo všetkých piatich riadkoch?} \\
            }
        }
    }
     \\ \\

    \theutterance \stepcounter{utterance}  
    & & & \multicolumn{2}{p{0.3\linewidth}}{
        \cellcolor[rgb]{0.9,0.9,0.9}{
            \makecell[{{p{\linewidth}}}]{
                \texttt{\tiny{[GM$|$GM]}}
                \texttt{continue} \\
            }
        }
    }
    & & \\ \\

    \theutterance \stepcounter{utterance}  
    & & \multicolumn{4}{p{0.6\linewidth}}{
        \cellcolor[rgb]{0.9,0.9,0.9}{
            \makecell[{{p{\linewidth}}}]{
                \texttt{\tiny{[P1$\langle$GM]}}
                \texttt{ODPOVEĎ: Áno, potvrdzujem, že v treťom riadku sú stĺpce 4 a 5 prázdne.} \\
\texttt{OTÁZKA: V prvom stĺpci sú podľa vášho popisu vo všetkých riadkoch znaky X. Potvrdzujete, že v prvom stĺpci sú naozaj znaky X vo všetkých piatich riadkoch?Svoju odpoveď začnite s "ODPOVEĎ:" a nepridávajte nič iné.} \\
            }
        }
    }
    & \\ \\

    \theutterance \stepcounter{utterance}  
    & \multicolumn{4}{p{0.6\linewidth}}{
        \cellcolor[rgb]{0.9,0.9,0.9}{
            \makecell[{{p{\linewidth}}}]{
                \texttt{\tiny{[P1$\rangle$GM]}}
                \texttt{ODPOVEĎ: Áno, potvrdzujem, že v prvom stĺpci sú znaky X vo všetkých piatich riadkoch.} \\
            }
        }
    }
    & & \\ \\

    \theutterance \stepcounter{utterance}  
    & & & \multicolumn{2}{p{0.3\linewidth}}{
        \cellcolor[rgb]{0.9,0.9,0.9}{
            \makecell[{{p{\linewidth}}}]{
                \texttt{\tiny{[GM$|$GM]}}
                \texttt{continue} \\
            }
        }
    }
    & & \\ \\

    \theutterance \stepcounter{utterance}  
    & & \multicolumn{4}{p{0.6\linewidth}}{
        \cellcolor[rgb]{0.9,0.9,0.9}{
            \makecell[{{p{\linewidth}}}]{
                \texttt{\tiny{[P1$\langle$GM]}}
                \texttt{Teraz položte otázku, aby ste zistili nové aspekty mojej mriežky, ktoré sa môžu líšiť od vašej mriežky. Začnite s "OTÁZKA:" a nepridávajte nič iné.} \\
            }
        }
    }
    & \\ \\

    \theutterance \stepcounter{utterance}  
    & \multicolumn{4}{p{0.6\linewidth}}{
        \cellcolor[rgb]{0.9,0.9,0.9}{
            \makecell[{{p{\linewidth}}}]{
                \texttt{\tiny{[P1$\rangle$GM]}}
                \texttt{OTÁZKA: V druhom riadku je podľa vášho popisu len jedno X na začiatku. Potvrdzujete, že v druhom riadku sú stĺpce 2 až 5 naozaj prázdne?} \\
            }
        }
    }
    & & \\ \\

    \theutterance \stepcounter{utterance}  
    & & & \multicolumn{2}{p{0.3\linewidth}}{
        \cellcolor[rgb]{0.9,0.9,0.9}{
            \makecell[{{p{\linewidth}}}]{
                \texttt{\tiny{[GM$|$GM]}}
                \texttt{continue} \\
            }
        }
    }
    & & \\ \\

    \theutterance \stepcounter{utterance}  
    & & \multicolumn{4}{p{0.6\linewidth}}{
        \cellcolor[rgb]{0.9,0.9,0.9}{
            \makecell[{{p{\linewidth}}}]{
                \texttt{\tiny{[GM$\rangle$P2]}}
                \texttt{ODPOVEĎ: Áno, potvrdzujem, že v prvom stĺpci sú znaky X vo všetkých piatich riadkoch.} \\
\texttt{OTÁZKA: V druhom riadku je podľa vášho popisu len jedno X na začiatku. Potvrdzujete, že v druhom riadku sú stĺpce 2 až 5 naozaj prázdne?Svoju odpoveď začnite s "ODPOVEĎ:" a nepridávajte nič iné.} \\
            }
        }
    }
    & \\ \\

    \theutterance \stepcounter{utterance}  
    & & & \multicolumn{4}{p{0.6\linewidth}}{
        \cellcolor[rgb]{0.9,0.9,0.9}{
            \makecell[{{p{\linewidth}}}]{
                \texttt{\tiny{[GM$\langle$P2]}}
                \texttt{ODPOVEĎ: Áno, potvrdzujem, že v druhom riadku sú stĺpce 2 až 5 prázdne.} \\
            }
        }
    }
     \\ \\

    \theutterance \stepcounter{utterance}  
    & & & \multicolumn{2}{p{0.3\linewidth}}{
        \cellcolor[rgb]{0.9,0.9,0.9}{
            \makecell[{{p{\linewidth}}}]{
                \texttt{\tiny{[GM$|$GM]}}
                \texttt{continue} \\
            }
        }
    }
    & & \\ \\

    \theutterance \stepcounter{utterance}  
    & & \multicolumn{4}{p{0.6\linewidth}}{
        \cellcolor[rgb]{0.9,0.9,0.9}{
            \makecell[{{p{\linewidth}}}]{
                \texttt{\tiny{[GM$\rangle$P2]}}
                \texttt{Teraz sa rozhodnite. Čo si myslíte: je vaša mriežka a druhá opísaná mriežka rovnaká? Napíšte "ROZHODNUTIE: rovnaká mriežka" ak si myslíte, že ide o rovnakú mriežku, alebo "ROZHODNUTIE: rôzne mriežky" ak si myslíte, že ide o rôzne mriežky. Nepridávajte nič iné.} \\
            }
        }
    }
    & \\ \\

    \theutterance \stepcounter{utterance}  
    & & & \multicolumn{4}{p{0.6\linewidth}}{
        \cellcolor[rgb]{0.9,0.9,0.9}{
            \makecell[{{p{\linewidth}}}]{
                \texttt{\tiny{[GM$\langle$P2]}}
                \texttt{ROZHODNUTIE: rôzne mriežky} \\
            }
        }
    }
     \\ \\

    \theutterance \stepcounter{utterance}  
    & & & \multicolumn{2}{p{0.3\linewidth}}{
        \cellcolor[rgb]{0.9,0.9,0.9}{
            \makecell[{{p{\linewidth}}}]{
                \texttt{\tiny{[GM$|$GM]}}
                \texttt{continue} \\
            }
        }
    }
    & & \\ \\

    \theutterance \stepcounter{utterance}  
    & & & \multicolumn{2}{p{0.3\linewidth}}{
        \cellcolor[rgb]{0.9,0.9,0.9}{
            \makecell[{{p{\linewidth}}}]{
                \texttt{\tiny{[GM$|$GM]}}
                \texttt{success} \\
            }
        }
    }
    & & \\ \\

    \theutterance \stepcounter{utterance}  
    & & \multicolumn{4}{p{0.6\linewidth}}{
        \cellcolor[rgb]{0.9,0.9,0.9}{
            \makecell[{{p{\linewidth}}}]{
                \texttt{\tiny{[P1$\langle$GM]}}
                \texttt{ODPOVEĎ: Áno, potvrdzujem, že v druhom riadku sú stĺpce 2 až 5 prázdne.} \\
\texttt{Teraz sa rozhodnite. Čo si myslíte: je vaša mriežka a druhá opísaná mriežka rovnaká? Napíšte "ROZHODNUTIE: rovnaká mriežka" ak si myslíte, že ide o rovnakú mriežku, alebo "ROZHODNUTIE: rôzne mriežky" ak si myslíte, že ide o rôzne mriežky. Nepridávajte nič iné.} \\
            }
        }
    }
    & \\ \\

    \theutterance \stepcounter{utterance}  
    & \multicolumn{4}{p{0.6\linewidth}}{
        \cellcolor[rgb]{0.9,0.9,0.9}{
            \makecell[{{p{\linewidth}}}]{
                \texttt{\tiny{[P1$\rangle$GM]}}
                \texttt{ROZHODNUTIE: rôzne mriežky} \\
            }
        }
    }
    & & \\ \\

    \theutterance \stepcounter{utterance}  
    & & & \multicolumn{2}{p{0.3\linewidth}}{
        \cellcolor[rgb]{0.9,0.9,0.9}{
            \makecell[{{p{\linewidth}}}]{
                \texttt{\tiny{[GM$|$GM]}}
                \texttt{continue} \\
            }
        }
    }
    & & \\ \\

    \theutterance \stepcounter{utterance}  
    & & & \multicolumn{2}{p{0.3\linewidth}}{
        \cellcolor[rgb]{0.9,0.9,0.9}{
            \makecell[{{p{\linewidth}}}]{
                \texttt{\tiny{[GM$|$GM]}}
                \texttt{success} \\
            }
        }
    }
    & & \\ \\

\end{supertabular}
}

\end{document}
